% ============================================================
%  10-release2.tex  (fusionné dans le chapitre 3)
%  Contient : Sprint 3 (Dimensioning Service + IA + PDF)
% ============================================================

% ══════════════════════════════════════════════════════════
\section{Sprint 3~: Service de Dimensionnement Solaire}
% ══════════════════════════════════════════════════════════

\subsection{Backlog du sprint}

\begin{kanbanbox}[breakable]{Sprint 3 --- Dimensioning Service (Port 8083)}
\centering\small
\begin{tabular}{|C{1cm}|L{4.5cm}|L{4cm}|C{2cm}|C{1.8cm}|}
  \hline
  \rowcolor{figblue!20}
  \textbf{ID Jira} & \textbf{Fonctionnalité} & \textbf{Tâches techniques} &
  \textbf{Priorité} & \textbf{Statut} \\
  \hline
  SOLAR-11 & Collecte données conso. & Formulaire multi-appareils, kWh/j calculé & Haute & \ok~Fait \\
  \hline
  SOLAR-07 & Intégration PVGIS & GET PVGIS API, GHI, kWh/kWc/j & Haute & \ok~Fait \\
  \hline
  SOLAR-13 & Calcul Night Panel & kWc, nb panneaux, surface, coût & Haute & \ok~Fait \\
  \hline
  SOLAR-10 & Assistant IA & Ollama/TinyLLaMA, recommandations & Moyenne & \ok~Fait \\
  \hline
  SOLAR-12 & Rapport PDF & iText 8, mise en page dynamique & Haute & \ok~Fait \\
  \hline
  SOLAR-12b & Historique calculs & GET /calculations/history, pagination & Moyenne & \ok~Fait \\
  \hline
  SOLAR-07b & Visualisation graphes & Chart.js, courbe mensuelle rayonnement & Basse & \ok~Fait \\
  \hline
\end{tabular}
\end{kanbanbox}

\subsection{Analyse du sprint 3}

Pour mieux comprendre ce troisième sprint, nous allons présenter le diagramme
de cas d'utilisation du sprint~3, décrire les différents cas d'utilisation,
et détailler les fonctionnalités mentionnées dans ce sprint.

\subsubsection{Diagramme de cas d'utilisation global relatif au sprint 3}

Ce diagramme de cas d'utilisation illustre les différentes fonctionnalités
prises en charge lors du sprint~3, accessibles à l'\textbf{Installateur},
qui réalise les études techniques de dimensionnement, exploite les
recommandations IA et fournit le rapport PDF au client.

\begin{figure}[H]
  \centering
  \IfFileExists{uc-s3-global.png}{%
    \includegraphics[width=0.88\textwidth]{uc-s3-global.png}%
  }{%
    \fbox{\parbox[c][8cm][c]{0.92\textwidth}{\centering Diagramme de cas d'utilisation global de Sprint 3}}%
  }
  \caption{Diagramme de cas d'utilisation global de Sprint~3}
  \label{fig:uc-s3}
\end{figure}

\subsubsection{Vue détaillée des cas d'utilisation développés durant le sprint 3}

Ensuite, nous présentons les diagrammes de cas d'utilisation détaillés
correspondant aux fonctionnalités principales. Ces diagrammes permettent
de mieux illustrer les interactions entre les acteurs et le système.

% ──────────────────────────────────────────────
\paragraph{Raffinement de cas d'utilisation «~Lancer le dimensionnement~»}

\begin{figure}[H]
  \centering
  \IfFileExists{uc-installateur-dimensionnement.png}{%
    \includegraphics[width=0.85\textwidth]{uc-installateur-dimensionnement.png}%
  }{%
    \fbox{\parbox[c][6cm][c]{0.85\textwidth}{\centering Diagramme de cas d'utilisation -- Lancer le dimensionnement}}%
  }
  \caption{Diagramme de cas d'utilisation pour «~Lancer le dimensionnement~»}
  \label{fig:uc-dimensionnement}
\end{figure}

\subparagraph{Description textuelle du cas d'utilisation «~Saisir les données de consommation~»}

\begin{table}[H]
  \caption{Description textuelle du cas d'utilisation «~Saisir les données de consommation~»}
  \label{tab:cu-saisir-conso}
  \centering
  \begin{tabularx}{\textwidth}{|L{3.8cm}|X|}
    \hline
    \rowcolor{figblue!10}
    \textbf{Cas d'utilisation} & Saisir les données de consommation. \\
    \hline
    \textbf{Acteur initiateur} & Installateur \\
    \hline
    \textbf{But du cas} & Permettre à l'installateur de renseigner la consommation électrique du client (par appareil ou agrégée en kWh/j) afin de préparer le dimensionnement. \\
    \hline
    \textbf{Précondition} & L'installateur est authentifié et un projet existe pour le client. \\
    \hline
    \textbf{Description du scénario} &
    \textbf{Scénario principal :}\newline
    1. L'installateur accède à la fiche projet du client.\newline
    2. Il clique sur «~Nouveau dimensionnement~».\newline
    3. Il choisit le mode de saisie~: par appareil (puissance, heures, quantité) ou total kWh/j.\newline
    4. Il saisit les valeurs et valide le formulaire.\newline
    5. Le frontend agrège les consommations par appareil pour calculer le total.\newline
    6. Les données sont stockées temporairement en attendant le lancement du calcul.\\
    \hline
    \textbf{Scénarios alternatifs} &
    \textit{Scénario alternatif 1 -- Consommation à zéro :}\newline
    \hspace*{6pt}4.1.~Le système détecte un total nul.\newline
    \hspace*{6pt}4.2.~Le système affiche : «~La consommation doit être supérieure à zéro~».\\
    \hline
  \end{tabularx}
\end{table}

\subparagraph{Description textuelle du cas d'utilisation «~Lancer le calcul de dimensionnement~»}

\begin{table}[H]
  \caption{Description textuelle du cas d'utilisation «~Lancer le calcul de dimensionnement~»}
  \label{tab:cu-lancer-calcul}
  \centering
  \begin{tabularx}{\textwidth}{|L{3.8cm}|X|}
    \hline
    \rowcolor{figblue!10}
    \textbf{Cas d'utilisation} & Lancer le calcul de dimensionnement. \\
    \hline
    \textbf{Acteur initiateur} & Installateur \\
    \hline
    \textbf{But du cas} & Permettre à l'installateur de calculer la puissance crête requise, le nombre de panneaux, la surface nécessaire et le coût estimé à partir de la consommation et de la localisation, en exploitant les données solaires PVGIS. \\
    \hline
    \textbf{Précondition} & Les données de consommation sont saisies, la localisation (latitude, longitude) est renseignée et l'API PVGIS est accessible. \\
    \hline
    \textbf{Description du scénario} &
    \textbf{Scénario principal :}\newline
    1. L'installateur valide le formulaire complet (consommation + localisation).\newline
    2. Le frontend envoie \texttt{POST /api/calculations} avec le JWT.\newline
    3. Le Dimensioning Service reçoit les paramètres.\newline
    4. \texttt{PvgisClient} interroge l'API PVGIS (latitude, longitude).\newline
    5. Récupération du GHI mensuel et du ratio kWh/kWc/jour.\newline
    6. \texttt{NightPanelCalculator} applique le modèle~: kWc, nb panneaux, surface, coût estimé.\newline
    7. Le résultat est persisté dans la table \texttt{dimensionings}.\newline
    8. Le Project Service est notifié pour passer le projet en \texttt{IN\_PROGRESS}.\newline
    9. La réponse complète est retournée au frontend.\\
    \hline
    \textbf{Scénarios alternatifs} &
    \textit{Scénario alternatif 1 -- PVGIS indisponible :}\newline
    \hspace*{6pt}4.1.~Le système renvoie HTTP~503 «~Service météo indisponible~».\newline[2pt]
    \textit{Scénario alternatif 2 -- Localisation hors couverture :}\newline
    \hspace*{6pt}5.1.~Le système indique que la zone n'est pas couverte par PVGIS.\\
    \hline
  \end{tabularx}
\end{table}

\subparagraph{Description textuelle du cas d'utilisation «~Consulter le résultat Night Panel~»}

\begin{table}[H]
  \caption{Description textuelle du cas d'utilisation «~Consulter le résultat Night Panel~»}
  \label{tab:cu-resultat-np}
  \centering
  \begin{tabularx}{\textwidth}{|L{3.8cm}|X|}
    \hline
    \rowcolor{figblue!10}
    \textbf{Cas d'utilisation} & Consulter le résultat du dimensionnement. \\
    \hline
    \textbf{Acteur initiateur} & Installateur \\
    \hline
    \textbf{But du cas} & Permettre à l'installateur de visualiser de façon synthétique les résultats du calcul : puissance crête, nombre de panneaux, surface, coût estimé, GHI mensuel et productible. \\
    \hline
    \textbf{Précondition} & Un calcul de dimensionnement a été effectué pour ce projet. \\
    \hline
    \textbf{Description du scénario} &
    \textbf{Scénario principal :}\newline
    1. L'installateur accède à la page «~Résultat du dimensionnement~».\newline
    2. Le frontend envoie \texttt{GET /api/calculations/\{id\}}.\newline
    3. Le système retourne les indicateurs clés et les données mensuelles.\newline
    4. L'écran affiche~: cartes KPI (kWc, panneaux, surface, coût), graphique mensuel du productible, carte de localisation.\newline
    5. L'installateur peut comparer avec d'autres calculs depuis l'historique.\\
    \hline
    \textbf{Scénarios alternatifs} &
    \textit{Scénario alternatif 1 -- Calcul inexistant :}\newline
    \hspace*{6pt}3.1.~L'identifiant est invalide.\newline
    \hspace*{6pt}3.2.~Le système renvoie HTTP~404.\\
    \hline
  \end{tabularx}
\end{table}

\subparagraph{Description textuelle du cas d'utilisation «~Consulter l'historique des calculs~»}

\begin{table}[H]
  \caption{Description textuelle du cas d'utilisation «~Consulter l'historique des calculs~»}
  \label{tab:cu-historique-calculs}
  \centering
  \begin{tabularx}{\textwidth}{|L{3.8cm}|X|}
    \hline
    \rowcolor{figblue!10}
    \textbf{Cas d'utilisation} & Consulter l'historique des calculs. \\
    \hline
    \textbf{Acteur initiateur} & Installateur \\
    \hline
    \textbf{But du cas} & Permettre à l'installateur de retrouver l'ensemble des calculs réalisés pour un projet et de comparer plusieurs simulations. \\
    \hline
    \textbf{Précondition} & L'installateur est authentifié et au moins un calcul a été effectué pour le projet. \\
    \hline
    \textbf{Description du scénario} &
    \textbf{Scénario principal :}\newline
    1. L'installateur accède à l'onglet «~Historique~» d'un projet.\newline
    2. Le frontend envoie \texttt{GET /api/calculations?projectId=\{id\}}.\newline
    3. Le système retourne la liste des calculs triés par date décroissante.\newline
    4. L'écran affiche un tableau avec date, kWc, nb panneaux, coût.\newline
    5. L'installateur peut sélectionner un ou plusieurs calculs pour les visualiser ou les comparer.\\
    \hline
    \textbf{Scénarios alternatifs} &
    \textit{Scénario alternatif 1 -- Historique vide :}\newline
    \hspace*{6pt}3.1.~Le projet n'a aucun calcul.\newline
    \hspace*{6pt}3.2.~Le système affiche : «~Aucun calcul effectué~».\\
    \hline
  \end{tabularx}
\end{table}

% ──────────────────────────────────────────────
\paragraph{Raffinement de cas d'utilisation «~Obtenir les recommandations IA~»}

\begin{figure}[H]
  \centering
  \IfFileExists{uc-installateur-recommendations.png}{%
    \includegraphics[width=0.85\textwidth]{uc-installateur-recommendations.png}%
  }{%
    \fbox{\parbox[c][6cm][c]{0.85\textwidth}{\centering Diagramme de cas d'utilisation -- Obtenir les recommandations IA}}%
  }
  \caption{Diagramme de cas d'utilisation pour «~Obtenir les recommandations IA~»}
  \label{fig:uc-recommendations}
\end{figure}

\subparagraph{Description textuelle du cas d'utilisation «~Demander des recommandations IA~»}

\begin{table}[H]
  \caption{Description textuelle du cas d'utilisation «~Demander des recommandations IA~»}
  \label{tab:cu-demander-reco}
  \centering
  \begin{tabularx}{\textwidth}{|L{3.8cm}|X|}
    \hline
    \rowcolor{figblue!10}
    \textbf{Cas d'utilisation} & Demander des recommandations IA. \\
    \hline
    \textbf{Acteur initiateur} & Installateur \\
    \hline
    \textbf{But du cas} & Permettre à l'installateur d'obtenir, à partir du
      résultat du dimensionnement, un \textbf{kit complet validé}~: onduleur
      Solax/Sungrow du catalogue société, câbles photovoltaïques tunisiens et
      disjoncteurs Schneider AC/DC. Le kit est accompagné d'un \textbf{score de
      compatibilité 0--100}, d'alertes éventuelles, d'arguments à présenter au
      client et d'une \textit{checklist} terrain. La pertinence est assurée par
      un pipeline RAG hybride~: sélection déterministe + recherche sémantique
      \texttt{pgvector} + génération Ollama. \\
    \hline
    \textbf{Précondition} & Le résultat du dimensionnement est disponible.
      La base \texttt{knowledge\_chunks} (pgvector) a été indexée au démarrage
      du service. \\
    \hline
    \textbf{Description du scénario} &
    \textbf{Scénario principal :}\newline
    1. L'installateur clique sur «~Recommandations IA~» depuis la page résultat.\newline
    2. Le frontend envoie \texttt{POST /api/dimensioning/\{id\}/ai/regenerate}.\newline
    3. Le \texttt{DecisionSupportService} sélectionne de manière déterministe
       l'onduleur (Solax/Sungrow, monophasé ou triphasé selon STEG), les
       câbles (NFC~33-209) et les disjoncteurs Schneider via
       \texttt{EquipmentSelectionService}.\newline
    4. \texttt{CompatibilityRuleEngine} calcule un score 0--100 et un verdict
       (\texttt{OK}, \texttt{ATTENTION}, \texttt{NON\_COMPATIBLE}).\newline
    5. \texttt{EmbeddingService} (Ollama \texttt{nomic-embed-text}) produit le
       vecteur de la requête et \texttt{KnowledgeChunkRepository} interroge
       \texttt{pgvector} (\texttt{ORDER BY embedding <=> \$1 LIMIT 5}) pour
       récupérer les passages pertinents (catalogue, normes, STEG/ANME).\newline
    6. Un prompt structuré est envoyé à Ollama qui retourne un JSON
       (alertes, arguments client, \textit{checklist}, narratif).\newline
    7. Le service fusionne le verdict déterministe avec l'enrichissement LLM
       et retourne un \texttt{InstallerRecommendationDto}.\newline
    8. Le frontend affiche le kit, le score, les alertes et les arguments. \\
    \hline
    \textbf{Scénarios alternatifs} &
    \textit{Scénario alternatif 1 -- Ollama indisponible :}\newline
    \hspace*{6pt}5.1.~Le service applique le \textbf{mode \textit{fallback}}~:
      le kit déterministe et le score sont renvoyés avec
      \texttt{fallback=true} et un narratif généré localement.\newline[2pt]
    \textit{Scénario alternatif 2 -- Catalogue insuffisant :}\newline
    \hspace*{6pt}3.1.~Si aucun onduleur ne couvre la puissance demandée,
      le plus puissant disponible est proposé et une alerte
      «~kit sous-dimensionné~» est ajoutée. \\
    \hline
  \end{tabularx}
\end{table}

\subparagraph{Description textuelle du cas d'utilisation «~Régénérer les recommandations~»}

\begin{table}[H]
  \caption{Description textuelle du cas d'utilisation «~Régénérer les recommandations~»}
  \label{tab:cu-regenerer-reco}
  \centering
  \begin{tabularx}{\textwidth}{|L{3.8cm}|X|}
    \hline
    \rowcolor{figblue!10}
    \textbf{Cas d'utilisation} & Régénérer les recommandations IA. \\
    \hline
    \textbf{Acteur initiateur} & Installateur \\
    \hline
    \textbf{But du cas} & Permettre à l'installateur de relancer le pipeline
      RAG (\textit{retrieval} \texttt{pgvector} + génération Ollama) s'il a
      modifié les paramètres du calcul ou s'il souhaite obtenir une nouvelle
      formulation du narratif et des arguments client. \\
    \hline
    \textbf{Précondition} & Un dimensionnement existe pour ce projet. \\
    \hline
    \textbf{Description du scénario} &
    \textbf{Scénario principal :}\newline
    1. L'installateur clique sur «~Régénérer~» dans la zone IA.\newline
    2. Le frontend envoie \texttt{POST /api/dimensioning/\{id\}/ai/regenerate}.\newline
    3. La sélection déterministe (onduleur, câbles, disjoncteurs) est
       recalculée afin de refléter les éventuelles modifications.\newline
    4. Une nouvelle requête vectorielle est exécutée sur \texttt{knowledge\_chunks}.\newline
    5. Ollama produit un nouveau JSON (alertes, arguments, \textit{checklist},
       narratif).\newline
    6. Le \texttt{InstallerRecommendationDto} ainsi qu'\texttt{aiRecommendation}
       sont persistés et remplacent les valeurs précédentes. \\
    \hline
    \textbf{Scénarios alternatifs} &
    \textit{Scénario alternatif 1 -- Quota dépassé :}\newline
    \hspace*{6pt}1.1.~L'installateur a régénéré plus de 3~fois en moins de 5~minutes.\newline
    \hspace*{6pt}1.2.~Le système refuse l'opération avec un message d'attente.\\
    \hline
  \end{tabularx}
\end{table}

% ──────────────────────────────────────────────
\paragraph{Raffinement de cas d'utilisation «~Télécharger le rapport PDF~»}

\begin{figure}[H]
  \centering
  \IfFileExists{uc-installateur-rapport-pdf.png}{%
    \includegraphics[width=0.85\textwidth]{uc-installateur-rapport-pdf.png}%
  }{%
    \fbox{\parbox[c][6cm][c]{0.85\textwidth}{\centering Diagramme de cas d'utilisation -- Télécharger le rapport PDF}}%
  }
  \caption{Diagramme de cas d'utilisation pour «~Télécharger le rapport PDF~»}
  \label{fig:uc-rapport-pdf}
\end{figure}

\subparagraph{Description textuelle du cas d'utilisation «~Télécharger le rapport PDF~»}

\begin{table}[H]
  \caption{Description textuelle du cas d'utilisation «~Télécharger le rapport PDF~»}
  \label{tab:cu-telecharger-pdf}
  \centering
  \begin{tabularx}{\textwidth}{|L{3.8cm}|X|}
    \hline
    \rowcolor{figblue!10}
    \textbf{Cas d'utilisation} & Télécharger le rapport technique PDF. \\
    \hline
    \textbf{Acteur initiateur} & Installateur \\
    \hline
    \textbf{But du cas} & Permettre à l'installateur de générer et télécharger un rapport PDF complet (logo entreprise, informations client, données PVGIS, résultats Night Panel, graphiques mensuels, recommandations IA) qu'il pourra remettre au client. \\
    \hline
    \textbf{Précondition} & Un calcul de dimensionnement a été effectué pour le projet. \\
    \hline
    \textbf{Description du scénario} &
    \textbf{Scénario principal :}\newline
    1. L'installateur clique sur «~Télécharger le rapport~» depuis la page résultat.\newline
    2. Le frontend envoie \texttt{GET /api/calculations/\{id\}/pdf}.\newline
    3. Le service récupère les données du calcul, du client et du projet.\newline
    4. \texttt{PdfGenerationService} construit le document via iText~8 (logo, infos client, données PVGIS, résultats Night Panel, graphique mensuel, recommandations IA).\newline
    5. Le fichier est retourné avec \texttt{Content-Type: application/pdf}.\newline
    6. Le navigateur déclenche le téléchargement \texttt{solarease-report-\{id\}.pdf}.\\
    \hline
    \textbf{Scénarios alternatifs} &
    \textit{Scénario alternatif 1 -- Calcul inexistant :}\newline
    \hspace*{6pt}3.1.~Le système renvoie HTTP~404.\newline[2pt]
    \textit{Scénario alternatif 2 -- Erreur génération :}\newline
    \hspace*{6pt}4.1.~iText échoue (police manquante, image non trouvée).\newline
    \hspace*{6pt}4.2.~Le système renvoie HTTP~500 avec un message diagnostic.\\
    \hline
  \end{tabularx}
\end{table}

\subparagraph{Description textuelle du cas d'utilisation «~Prévisualiser le rapport~»}

\begin{table}[H]
  \caption{Description textuelle du cas d'utilisation «~Prévisualiser le rapport~»}
  \label{tab:cu-previsualiser-pdf}
  \centering
  \begin{tabularx}{\textwidth}{|L{3.8cm}|X|}
    \hline
    \rowcolor{figblue!10}
    \textbf{Cas d'utilisation} & Prévisualiser le rapport avant téléchargement. \\
    \hline
    \textbf{Acteur initiateur} & Installateur \\
    \hline
    \textbf{But du cas} & Permettre à l'installateur de visualiser le rapport directement dans le navigateur avant de le télécharger ou de le partager. \\
    \hline
    \textbf{Précondition} & Le rapport est généré pour un calcul existant. \\
    \hline
    \textbf{Description du scénario} &
    \textbf{Scénario principal :}\newline
    1. L'installateur clique sur «~Prévisualiser~» depuis la page résultat.\newline
    2. Le frontend ouvre le PDF dans un \texttt{iframe} ou une modale.\newline
    3. L'installateur navigue dans les pages du rapport.\newline
    4. Il peut alors cliquer sur «~Télécharger~» ou «~Partager~».\\
    \hline
    \textbf{Scénarios alternatifs} &
    \textit{Scénario alternatif 1 -- Navigateur non compatible :}\newline
    \hspace*{6pt}2.1.~Le système propose un téléchargement direct.\\
    \hline
  \end{tabularx}
\end{table}

\subsection{Conception du sprint 3}

L'analyse conceptuelle de ce troisième sprint comprendra la description des
diagrammes de séquence liés aux fonctionnalités majeures du sprint~3. Elle
se conclura par l'élaboration du diagramme de classes, synthétisant
l'ensemble des éléments étudiés.

\subsubsection{Diagrammes de séquence}

Dans cette phase, nous présentons les diagrammes de séquence correspondant
aux principales fonctionnalités développées lors du sprint~3. Ces diagrammes
illustrent les interactions entre les différents composants du système pour
chaque fonctionnalité ciblée.

% ──────────────────────────────────────────────
\paragraph{Diagramme de séquence «~Lancer le dimensionnement~»}

La figure ci-après présente le scénario nominal d'un calcul de dimensionnement~:
saisie de la consommation, appel à l'API PVGIS pour récupérer les données
solaires, application du modèle Night Panel, persistance du résultat.

\begin{figure}[H]
  \centering
  \IfFileExists{seq-dimensionnement.png}{%
    \includegraphics[width=0.95\textwidth,height=0.7\textheight,keepaspectratio]{seq-dimensionnement.png}%
  }{%
    \fbox{\parbox[c][7cm][c]{0.92\textwidth}{\centering Diagramme de séquence «~Lancer le dimensionnement~»}}%
  }
  \caption{Diagramme de séquence «~Lancer le dimensionnement~»}
  \label{fig:seq-s3}
\end{figure}

% ──────────────────────────────────────────────
\paragraph{Diagramme de séquence «~Obtenir les recommandations IA~»}

La figure ci-après illustre le pipeline RAG hybride~: sélection déterministe
du kit (\texttt{EquipmentSelectionService} + \texttt{CompatibilityRuleEngine}),
recherche sémantique dans la base vectorielle \texttt{pgvector}, génération
narrative par Ollama, puis restitution structurée
(\texttt{InstallerRecommendationDto}) au frontend.

\begin{figure}[H]
  \centering
  \IfFileExists{seq-recommandations-ia.png}{%
    \includegraphics[width=0.95\textwidth,height=0.7\textheight,keepaspectratio]{seq-recommandations-ia.png}%
  }{%
    \fbox{\parbox[c][7cm][c]{0.92\textwidth}{\centering Diagramme de séquence «~Obtenir les recommandations IA~»}}%
  }
  \caption{Diagramme de séquence «~Obtenir les recommandations IA~»}
  \label{fig:seq-reco-ia}
\end{figure}

% ──────────────────────────────────────────────
\paragraph{Diagramme de séquence «~Télécharger le rapport PDF~»}

La figure ci-après détaille la génération du rapport PDF via iText~8 à partir
des données du calcul, du client et des recommandations IA.

\begin{figure}[H]
  \centering
  \IfFileExists{seq-rapport-pdf.png}{%
    \includegraphics[width=0.95\textwidth,height=0.7\textheight,keepaspectratio]{seq-rapport-pdf.png}%
  }{%
    \fbox{\parbox[c][7cm][c]{0.92\textwidth}{\centering Diagramme de séquence «~Télécharger le rapport PDF~»}}%
  }
  \caption{Diagramme de séquence «~Télécharger le rapport PDF~»}
  \label{fig:seq-pdf}
\end{figure}

\subsubsection{Diagramme de classes du sprint 3}

Ce diagramme de classes global illustre l'architecture du projet, avec les
classes liées au sprint~3.

Il permet d'identifier clairement les entités utilisées lors de ce troisième
sprint.

\begin{figure}[H]
  \centering
  \IfFileExists{class-sprint3.png}{%
    \includegraphics[width=0.85\textwidth]{class-sprint3.png}%
  }{%
    \fbox{\parbox[c][8cm][c]{0.92\textwidth}{\centering Diagramme de classes global du sprint 3}}%
  }
  \caption{Diagramme de classes global du sprint~3}
  \label{fig:class-s3}
\end{figure}

\subsection{Réalisation du sprint 3}

\subsubsection{Modèle de calcul Night Panel}

\begin{lstlisting}[caption={NightPanelCalculator.java --- Logique de dimensionnement},
                   label=lst:nightpanel, language=Java]
@Component
public class NightPanelCalculator {

    // Parametres du modele Night Panel
    private static final double PANEL_EFFICIENCY   = 0.20; // 20%%
    private static final double SYSTEM_LOSSES      = 0.80; // pertes systeme
    private static final double PANEL_POWER_KWC    = 0.40; // 400 Wc / panneau
    private static final double COST_PER_KWC_TND   = 2_800; // TND/kWc

    public DimensionResult calculate(double dailyKwh, double ghi) {
        // 1. Puissance creete requise (kWc)
        double kWcRequired = dailyKwh / (ghi * SYSTEM_LOSSES);
        // 2. Nombre de panneaux (arrondi superieur)
        int nbPanels = (int) Math.ceil(kWcRequired / PANEL_POWER_KWC);
        // 3. Surface utilisee (m2)
        double surface = nbPanels * 1.6 * 0.99;      // 1.6 x 0.99 m
        // 4. Cout estime (TND)
        double cost = kWcRequired * COST_PER_KWC_TND;
        return new DimensionResult(kWcRequired, nbPanels, surface, cost);
    }
}
\end{lstlisting}

\begin{lstlisting}[caption={OllamaService.java --- Appel au modèle IA},
                   label=lst:ollama, language=Java]
@Service
public class OllamaService {

    public String generateRecommendations(DimensionResult result, double ghi) {
        String prompt = String.format(
            "Tu es un expert en energie solaire en Tunisie.\n" +
            "Puissance requise : %.2f kWc. Consommation journaliere : %.1f kWh.\n" +
            "Localisation GHI : %.2f kWh/m2/j. Nombre panneaux : %d.\n" +
            "Recommande : 1) Marques de panneaux, 2) Entretien, 3) Retour investissement.",
            result.getKwcRequired(), result.getDailyKwh(),
            ghi, result.getNbPanels());

        OllamaRequest req = new OllamaRequest("tinyllama", prompt, false);
        OllamaResponse res = restTemplate.postForObject(
            ollamaUrl + "/api/generate", req, OllamaResponse.class);
        return res != null ? res.getResponse() : "Service IA indisponible.";
    }
}
\end{lstlisting}

\subsubsection{Interfaces utilisateur réalisées}

\begin{figure}[H]
  \centering
  \fbox{\parbox{0.85\textwidth}{
    \centering\vspace{0.5cm}
    \textit{Formulaire de saisie des données de consommation}\\[4pt]
    \small Remplacer par : \texttt{\textbackslash includegraphics[width=0.85\textbackslash textwidth]\{images/consumption-form.png\}}\\[4pt]
    Liste d'appareils dynamique, puissance (W), heures/jour, quantité, total kWh/j auto-calculé
    \vspace{0.5cm}
  }}
  \caption{Formulaire de saisie de consommation -- Sprint 3}
  \label{fig:ui-conso}
\end{figure}

\begin{figure}[H]
  \centering
  \fbox{\parbox{0.85\textwidth}{
    \centering\vspace{0.5cm}
    \textit{Page de résultats du dimensionnement -- Modèle Night Panel}\\[4pt]
    \small Remplacer par : \texttt{\textbackslash includegraphics[width=0.85\textbackslash textwidth]\{images/dimensioning-result.png\}}\\[4pt]
    Puissance kWc, Nb panneaux, Surface (m²), Coût estimé (TND), Graphe mensuel PVGIS
    \vspace{0.5cm}
  }}
  \caption{Résultats de dimensionnement -- Sprint 3}
  \label{fig:ui-result}
\end{figure}

\begin{figure}[H]
  \centering
  \fbox{\parbox{0.85\textwidth}{
    \centering\vspace{0.5cm}
    \textit{Interface des recommandations IA (Ollama / TinyLLaMA)}\\[4pt]
    \small Remplacer par : \texttt{\textbackslash includegraphics[width=0.85\textbackslash textwidth]\{images/ai-recommendations.png\}}\\[4pt]
    Bulles de dialogue~: marques recommandées, entretien, ROI estimé, conseils personnalisés
    \vspace{0.5cm}
  }}
  \caption{Recommandations IA embarquée -- Sprint 3}
  \label{fig:ui-ia}
\end{figure}

\begin{figure}[H]
  \centering
  \fbox{\parbox{0.85\textwidth}{
    \centering\vspace{0.5cm}
    \textit{Aperçu du rapport PDF généré (iText 8)}\\[4pt]
    \small Remplacer par : \texttt{\textbackslash includegraphics[width=0.85\textbackslash textwidth]\{images/pdf-report.png\}}\\[4pt]
    Logo SolarEase, données client, tableau résultats, graphique ensoleillement, recommandations IA
    \vspace{0.5cm}
  }}
  \caption{Rapport PDF généré automatiquement -- Sprint 3}
  \label{fig:ui-pdf}
\end{figure}

\subsection{Environnement du sprint 3}

\begin{table}[H]
  \caption{Technologies et outils utilisés -- Sprint 3}
  \label{tab:tech-s3}
  \centering
  \begin{tabularx}{\textwidth}{|L{3cm}|L{3cm}|X|}
    \hline
    \rowcolor{figblue!20}
    \textbf{Technologie} & \textbf{Version} & \textbf{Rôle dans le sprint 3} \\
    \hline
    Spring Boot & 3.2 & Microservice de dimensionnement (Dimensioning Service) \\
    \hline
    PVGIS API & v5.2 & Données d'ensoleillement (GHI mensuel, kWh/kWc/j) \\
    \hline
    Ollama & 0.1.x & Runtime local pour le modèle TinyLLaMA (recommandations IA) \\
    \hline
    TinyLLaMA & 1.1B & Modèle de langage embarqué~; génération de recommandations \\
    \hline
    iText 8 & 8.0 & Génération de rapports PDF dynamiques et professionnels \\
    \hline
    Chart.js & 4.x & Visualisation des données PVGIS (rayonnement mensuel) \\
    \hline
    PostgreSQL & 16 & Stockage des calculs, historique et données de consommation \\
    \hline
  \end{tabularx}
\end{table}

\subsubsection{Tests et validation}

\begin{table}[H]
  \caption{Résultats des tests fonctionnels -- Sprint 3}
  \label{tab:tests-s3}
  \centering\small
  \begin{tabularx}{\textwidth}{|C{1.2cm}|L{3.5cm}|X|C{1.4cm}|}
    \hline
    \rowcolor{figblue!20}
    \textbf{ID} & \textbf{Cas de test} & \textbf{Résultat attendu} & \textbf{Statut} \\
    \hline
    T-17 & Dimensionnement valide & kWc, nb panneaux, surface, coût calculés & \ok~OK \\
    \hline
    T-18 & PVGIS disponible & GHI récupéré, données stockées & \ok~OK \\
    \hline
    T-19 & PVGIS hors service & Message d'erreur + données cache & \ok~OK \\
    \hline
    T-20 & Recommandations IA & Texte généré $<$ 3 secondes & \ok~OK \\
    \hline
    T-21 & Génération PDF & Fichier PDF valide, taille $<$ 2 Mo & \ok~OK \\
    \hline
    T-22 & Historique paginé & Liste des calculs + métadonnées & \ok~OK \\
    \hline
    T-23 & Résultat Night Panel & Cohérence kWc/nb panneaux & \ok~OK \\
    \hline
    T-24 & Accès non autorisé & HTTP 401 via Gateway & \ok~OK \\
    \hline
  \end{tabularx}
\end{table}

% ══════════════════════════════════════════════════════════
\subsection{Approche RAG hybride (pgvector + catalogue société)}
% ══════════════════════════════════════════════════════════

L'assistant IA de SolarEase ne se limite pas à un simple appel LLM.
Il combine~:

\begin{itemize}
  \item un \textbf{moteur déterministe} qui sélectionne le kit complet
    (onduleur Solax/Sungrow, câbles photovoltaïques tunisiens normalisés,
    disjoncteurs Schneider AC/DC) à partir du \textbf{catalogue de la
    société} stocké dans la table \texttt{equipments}~;
  \item un \textbf{moteur sémantique} (\textbf{RAG}, \textit{Retrieval-Augmented
    Generation}) qui interroge une base vectorielle \texttt{pgvector}
    (table \texttt{knowledge\_chunks}, embeddings de dimension~768)
    contenant la documentation produit, les normes câblage et les règles
    STEG/ANME~;
  \item un \textbf{LLM local} (Ollama) qui produit, à partir du kit et
    des passages retrouvés, une réponse \textbf{JSON structurée}
    (alertes, arguments client, \textit{checklist} terrain, narratif).
\end{itemize}

Cette architecture garantit que la \textbf{décision technique} (quel
onduleur, quel câble, quel disjoncteur) reste \textbf{reproductible et
auditable}, tandis que la \textbf{communication vers le client} bénéficie
de la richesse d'un LLM. En cas d'indisponibilité d'Ollama, le mode
\textit{fallback} renvoie le kit déterministe seul, sans rupture de service.

\subsubsection{Architecture du pipeline RAG hybride}

Le pipeline se décompose en quatre étapes~:

\begin{enumerate}
  \item \textbf{Sélection déterministe~:} \texttt{Equip\-mentSelectionService}
    interroge \texttt{equipments} pour choisir l'onduleur de plus petite
    puissance couvrant la capacité demandée et respectant la règle STEG
    monophasé/triphasé ($\leq$~5,5~kWc en monophasé, $>$~5,5~kWc en
    triphasé). Les sections de câbles DC/AC suivent la norme NFC~33-209
    en fonction de la puissance, et les disjoncteurs Schneider sont
    dimensionnés sur le courant de court-circuit majoré de 25~\%.
  \item \textbf{Scoring de compatibilité~:} \texttt{Compatibility\-RuleEngine}
    calcule un score 0--100 et un verdict (\texttt{OK},
    \texttt{ATTENTION}, \texttt{NON\_COMPATIBLE}) en pénalisant les
    écarts (sur-/sous-dimensionnement onduleur, mauvaise phase, câble
    sous-dimensionné, disjoncteur manquant). Les anomalies sont
    converties en \textbf{alertes textuelles}.
  \item \textbf{Retrieval sémantique~:} \texttt{Embed\-dingService}
    appelle Ollama (\texttt{nomic-embed-text}) pour transformer la
    requête en vecteur de dimension~768. \texttt{Knowledge\-Chunk\-Repository}
    interroge ensuite \texttt{pgvector} via l'opérateur de cosinus
    \texttt{embedding <=> \$1} et récupère les 5 \textit{chunks} les
    plus proches.
  \item \textbf{Generation~:} le \texttt{DecisionSupportService} construit
    un prompt incluant le kit, les alertes et les passages retrouvés,
    puis demande à Ollama un JSON conforme à un schéma précis
    (\texttt{alerts}, \texttt{clientArguments}, \texttt{terrainChecklist},
    \texttt{narrativeSummary}). Le service fusionne le résultat avec le
    verdict déterministe et renvoie un \texttt{InstallerRecommendationDto}.
\end{enumerate}

\subsubsection{Indexation de la base vectorielle}

Au démarrage du Dimensioning Service, \texttt{Knowledge\-IndexerService}
construit ou met à jour la table \texttt{knowledge\_chunks}~:

\begin{table}[H]
  \caption{Sources indexées dans la base vectorielle pgvector}
  \label{tab:kb-rag}
  \centering
  \begin{tabularx}{\textwidth}{|L{3.5cm}|X|}
    \hline
    \rowcolor{figblue!20}
    \textbf{Source} & \textbf{Contenu indexé (chunks)} \\
    \hline
    Catalogue société & Onduleurs Solax / Sungrow (plages de puissance,
      monophasé/triphasé, prix), panneaux et batteries de référence,
      disjoncteurs Schneider AC/DC (calibres, type \textit{photovoltaïque}). \\
    \hline
    Normes câblage & Câbles photovoltaïques tunisiens, sections par
      bracket de puissance (NFC~33-209), chutes de tension admissibles,
      températures d'usage. \\
    \hline
    Règles STEG / ANME & Plafonds d'autoconsommation, raccordement
      monophasé / triphasé, primes ANME, démarches administratives
      (autorisations, raccordement, comptage). \\
    \hline
  \end{tabularx}
\end{table}

\subsubsection{Schéma de la table \texttt{knowledge\_chunks}}

\begin{lstlisting}[caption={Schéma pgvector --- knowledge\_chunks},
                   label=lst:pgvector, language=SQL]
CREATE EXTENSION IF NOT EXISTS vector;

CREATE TABLE IF NOT EXISTS knowledge_chunks (
    id          BIGSERIAL PRIMARY KEY,
    source_type VARCHAR(40) NOT NULL,   -- EQUIPMENT | RULE | REGULATION
    source_id   VARCHAR(120),
    content     TEXT NOT NULL,
    metadata    JSONB,
    embedding   vector(768)             -- nomic-embed-text
);

CREATE INDEX IF NOT EXISTS knowledge_chunks_embedding_idx
    ON knowledge_chunks
    USING ivfflat (embedding vector_cosine_ops)
    WITH (lists = 100);
\end{lstlisting}

\subsubsection{Extrait de code --- Pipeline RAG hybride}

\begin{lstlisting}[caption={DecisionSupportService.java --- pipeline hybride pgvector},
                   label=lst:rag, language=Java]
public Result generate(DimensioningResponse dim,
                       Equipment usedPanel,
                       Equipment usedInverter) {
    double kw = dim.getInstallation().getTotalCapacityKw();

    // 1. Selection deterministe (catalogue societe)
    InstallerRecommendationDto.RecommendedKit kit = buildKit(kw, usedInverter);

    // 2. Scoring de compatibilite
    CompatibilityRuleEngine.Result eval = ruleEngine.evaluate(
            kw, kit.getInverter(), kit.getDcCable(), null, null);

    InstallerRecommendationDto dto = baseDto(eval, kit, kw, usedInverter);

    try {
        // 3. Retrieval semantique (pgvector)
        float[] q = embeddingService.embed(buildQuery(dim, kit));
        List<KnowledgeChunk> chunks =
                chunkRepository.searchSimilar(q, RETRIEVAL_TOPK);
        dto.getRagSources().addAll(toSourceLabels(chunks));

        // 4. Generation LLM + parsing JSON
        String json = ollamaService.generate(buildPrompt(dim, kit, dto, chunks));
        applyLlmAnswer(dto, json);
    } catch (Exception e) {
        log.warn("RAG augmentation skipped, deterministic kit only: {}",
                 e.getMessage());
        dto.setFallback(true);
        applyDeterministicNarrative(dto, dim.getInstallation(), dim.getFinancials());
    }
    return new Result(dto, dto.getNarrativeSummary());
}
\end{lstlisting}

\begin{conclusionbox}
Le sprint~3 concrétise la proposition de valeur centrale de SolarEase.
Le modèle Night Panel offre des calculs fiables validés contre des
références terrain. Le pipeline RAG \textbf{hybride} couple une sélection
déterministe sur le \textbf{catalogue société} (onduleurs Solax/Sungrow,
câbles tunisiens, disjoncteurs Schneider) à une recherche sémantique
\textbf{pgvector} et à un LLM local Ollama. L'installateur reçoit ainsi
un \textbf{kit prêt à commander}, un \textbf{score de compatibilité} et
des \textbf{arguments à présenter au client}, le tout avec un
\textit{fallback} déterministe garantissant la disponibilité même si le
modèle est hors-ligne. La génération PDF professionnelle permet enfin
à l'installateur de remettre un rapport complet au client final.
\end{conclusionbox}

\cleardoublepage
